\documentclass[11pt,letterpaper]{article}
\usepackage[technical-reference]{cornell-notes}
\usepackage{needspace}

\title{Clause 16: Library Introduction - Cornell Notes}
\author{}
\date{}
\setDocTitle{Clause 16: Library Introduction}
\setDocSubtitle{C++ 2024 Cornell Notes}
\setDocOwner{C++ 2024 Cornell Notes}
\setDocAuthor{}
\setDocDate{}
\setCornellCollection{C++ 2024 Cornell Notes}
\setCornellUnitType{Clause}
\setCornellUnitNumber{16}
\setCornellUnitTitle{Library Introduction}
\setCornellCompanionLabel{C++ Technical Reference}
\hypersetup{pdftitle={Clause 16: Library Introduction - Cornell Notes},pdfsubject={C++ technical study notes},pdfauthor={}}

\begin{document}
\maketitle
\begin{CornellReferenceOrientationBox}
\textbf{Purpose.} Defines how library clauses are organized and interpreted, then states requirements on programs and conforming implementations, including reserved names, user specializations, type requirements, moved-from states, reentrancy, and data-race avoidance.
\par\textbf{Role.} The interpretation and conformance framework for every library facility.
\par\textbf{Principal dependencies.} Uses the library-wide rules of Clause 16 together with the applicable core-language concepts and supporting library clauses.
\par\textbf{Primary audience.} programmers and library implementers.
\par\textbf{Major subject groups.} General; The C standard library; Method of description; Library-wide requirements.
\end{CornellReferenceOrientationBox}
\section{Learning Objectives}
\begin{itemize}[label=\textcolor{CornellReferenceTeal}{$\blacktriangleright$}]
\item Define the governing ideas of General and apply its stated contract at a call site.
\item Identify the governing ideas of The C standard library and apply its stated contract at a call site.
\item Distinguish the governing ideas of Method of description and apply its stated contract at a call site.
\item Explain the governing ideas of Library-wide requirements and apply its stated contract at a call site.
\item Trace the governing ideas of General and apply its stated contract at a call site.
\item Apply the governing ideas of Structure of each clause and apply its stated contract at a call site.
\item Analyze the governing ideas of Other conventions and apply its stated contract at a call site.
\item Evaluate the governing ideas of General and apply its stated contract at a call site.
\item Diagnose the governing ideas of Library contents and organization and apply its stated contract at a call site.
\item Compare the governing ideas of Using the library and apply its stated contract at a call site.
\item Verify the governing ideas of Requirements on types and expressions and apply its stated contract at a call site.
\item Relate the governing ideas of Constraints on programs and apply its stated contract at a call site.
\item Classify the governing ideas of Conforming implementations and apply its stated contract at a call site.
\item Justify the governing ideas of mandate and apply its stated contract at a call site.
\item Audit the governing ideas of precondition and apply its stated contract at a call site.
\item Summarize the governing ideas of exposition-only entity and apply its stated contract at a call site.
\end{itemize}
\section{Normative Structure Map}
\small\begin{longtable}{@{}>{\RaggedRight\arraybackslash}p{0.13\textwidth}>{\RaggedRight\arraybackslash}p{0.27\textwidth}>{\RaggedRight\arraybackslash}p{0.19\textwidth}>{\RaggedRight\arraybackslash}p{0.31\textwidth}@{}}
\toprule\textbf{Subclause} & \textbf{Subject} & \textbf{Rule category} & \textbf{Key dependencies / entities}\\\midrule\endfirsthead
\toprule\textbf{Subclause} & \textbf{Subject} & \textbf{Rule category} & \textbf{Key dependencies / entities}\\\midrule\endhead
\bottomrule\endfoot
16.1 & General & library semantics & Uses the library-wide rules of Clause 16 together with the applicable core-language concepts and supporting library clauses.\\
16.2 & The C standard library & library semantics & Uses the library-wide rules of Clause 16 together with the applicable core-language concepts and supporting library clauses.\\
16.3 & Method of description & library semantics & General, Structure of each clause, Other conventions\\
16.4 & Library-wide requirements & requirements & General, Library contents and organization, Using the library, Requirements on types and expressions, and related subclauses\\
\end{longtable}\normalsize
\begin{CornellReferenceNormativeBox}A heading maps the clause; it does not replace the detailed conditions. For each operation, read requirements, effects, result, exceptions, complexity, remarks, and notes with their stated normative force.\end{CornellReferenceNormativeBox}
\subsection*{Rule Classification Guide}
\small\begin{longtable}{@{}>{\RaggedRight\arraybackslash}p{0.25\textwidth}>{\RaggedRight\arraybackslash}p{0.67\textwidth}@{}}\toprule\textbf{Label or category} & \textbf{How to read it}\\\midrule\endfirsthead\toprule\textbf{Label or category} & \textbf{How to read it (continued)}\\\midrule\endhead\bottomrule\endfoot
Constraints & Conditions controlling whether a declaration, overload, or specialization participates in the relevant context.\\
Mandates & Requirements checked for the described declaration or specialization; failure has the stated well-formedness consequence.\\
Preconditions & Obligations the caller establishes before evaluation; violating one forfeits the operation's contract.\\
Effects / Returns / Postconditions & Separate the state change, produced result, and state guaranteed after successful completion.\\
Throws / Error conditions & State how failure is reported and which exceptional paths are permitted or required.\\
Complexity & Bounds or exact counts for the operations named by the specification.\\
Remarks / Recommended practice & Remarks refine applicability or participation; recommended practice guides implementations without silently becoming a program requirement.\\
Exposition-only & A device for describing semantics, not a public name on which a program can depend.\\
\end{longtable}\normalsize
\section{Cornell Cue-and-Notes}
\Needspace{8\baselineskip}
\subsection*{16.1 General}
\begin{longtable}{@{}>{\RaggedRight\arraybackslash\columncolor{CornellReferenceCueGray}}p{0.28\textwidth}>{\RaggedRight\arraybackslash}p{0.64\textwidth}@{}}
\rowcolor{CornellReferenceNavy}\color{white}\textbf{Cue / recall prompt} & \color{white}\textbf{Notes}\\\endfirsthead
\rowcolor{CornellReferenceNavy}\color{white}\textbf{Cue / recall prompt} & \color{white}\textbf{Notes (continued)}\\\endhead
\arrayrulecolor{CornellReferenceLineGray}
\textbf{What contract governs General?}\\[-1mm]{\scriptsize\CornellClauseRef{16.1}} & Frames the terminology, applicability, and relationships needed to interpret General; detailed rules remain in the following subclauses.\\\hline
\end{longtable}
\begin{CornellReferenceSummaryBox}General supplies the library semantics portion of this clause. The key review move is to connect its local wording to the clause-wide model, then classify each condition by normative force before relying on it.\end{CornellReferenceSummaryBox}
\Needspace{8\baselineskip}
\subsection*{16.2 The C standard library}
\begin{longtable}{@{}>{\RaggedRight\arraybackslash\columncolor{CornellReferenceCueGray}}p{0.28\textwidth}>{\RaggedRight\arraybackslash}p{0.64\textwidth}@{}}
\rowcolor{CornellReferenceNavy}\color{white}\textbf{Cue / recall prompt} & \color{white}\textbf{Notes}\\\endfirsthead
\rowcolor{CornellReferenceNavy}\color{white}\textbf{Cue / recall prompt} & \color{white}\textbf{Notes (continued)}\\\endhead
\arrayrulecolor{CornellReferenceLineGray}
\textbf{What contract governs The C standard library?}\\[-1mm]{\scriptsize\CornellClauseRef{16.2}} & Specifies the facility family The C standard library. Read its declarations together with mandates, preconditions, effects, results, exceptions, complexity, and lifetime rules.\\\hline
\end{longtable}
\begin{CornellReferenceSummaryBox}The C standard library supplies the library semantics portion of this clause. The key review move is to connect its local wording to the clause-wide model, then classify each condition by normative force before relying on it.\end{CornellReferenceSummaryBox}
\Needspace{8\baselineskip}
\subsection*{16.3 Method of description}
\begin{longtable}{@{}>{\RaggedRight\arraybackslash\columncolor{CornellReferenceCueGray}}p{0.28\textwidth}>{\RaggedRight\arraybackslash}p{0.64\textwidth}@{}}
\rowcolor{CornellReferenceNavy}\color{white}\textbf{Cue / recall prompt} & \color{white}\textbf{Notes}\\\endfirsthead
\rowcolor{CornellReferenceNavy}\color{white}\textbf{Cue / recall prompt} & \color{white}\textbf{Notes (continued)}\\\endhead
\arrayrulecolor{CornellReferenceLineGray}
\textbf{What contract governs General?}\\[-1mm]{\scriptsize\CornellClauseRef{16.3.1}} & Frames the terminology, applicability, and relationships needed to interpret General; detailed rules remain in the following subclauses.\\\hline
\textbf{What contract governs Structure of each clause?}\\[-1mm]{\scriptsize\CornellClauseRef{16.3.2}} & Specifies the facility family Structure of each clause. Read its declarations together with mandates, preconditions, effects, results, exceptions, complexity, and lifetime rules.\\\hline
\textbf{What contract governs Other conventions?}\\[-1mm]{\scriptsize\CornellClauseRef{16.3.3}} & Specifies the facility family Other conventions. Read its declarations together with mandates, preconditions, effects, results, exceptions, complexity, and lifetime rules.\\\hline
\end{longtable}
\begin{CornellReferenceSummaryBox}Method of description supplies the library semantics portion of this clause. The key review move is to connect its local wording to the clause-wide model, then classify each condition by normative force before relying on it.\end{CornellReferenceSummaryBox}
\Needspace{5\baselineskip}\paragraph{Detailed coverage ledger.}
\begin{description}[style=nextline,leftmargin=2.8cm,font=\normalfont\bfseries\color{CornellReferenceTeal}]
\item[16.3.2.1] \textbf{Elements.} Specifies the facility family Elements. Read its declarations together with mandates, preconditions, effects, results, exceptions, complexity, and lifetime rules.
\item[16.3.2.2] \textbf{Summary.} Specifies the facility family Summary. Read its declarations together with mandates, preconditions, effects, results, exceptions, complexity, and lifetime rules.
\item[16.3.2.3] \textbf{Requirements.} States the admissibility and semantic obligations for Requirements. Separate compile-time participation rules from caller preconditions and runtime effects.
\item[16.3.2.4] \textbf{Detailed specifications.} Specifies the facility family Detailed specifications. Read its declarations together with mandates, preconditions, effects, results, exceptions, complexity, and lifetime rules.
\item[16.3.2.5] \textbf{C library.} Specifies the facility family C library. Read its declarations together with mandates, preconditions, effects, results, exceptions, complexity, and lifetime rules.
\item[16.3.3.1] \textbf{General.} Frames the terminology, applicability, and relationships needed to interpret General; detailed rules remain in the following subclauses.
\item[16.3.3.2] \textbf{Exposition-only entities, etc..} Specifies the facility family Exposition-only entities, etc.. Read its declarations together with mandates, preconditions, effects, results, exceptions, complexity, and lifetime rules.
\item[16.3.3.3] \textbf{Type descriptions.} Specifies the facility family Type descriptions. Read its declarations together with mandates, preconditions, effects, results, exceptions, complexity, and lifetime rules.
\item[16.3.3.3.1] \textbf{General.} Frames the terminology, applicability, and relationships needed to interpret General; detailed rules remain in the following subclauses.
\item[16.3.3.3.2] \textbf{Enumerated types.} Specifies the facility family Enumerated types. Read its declarations together with mandates, preconditions, effects, results, exceptions, complexity, and lifetime rules.
\item[16.3.3.3.3] \textbf{Bitmask types.} Specifies the facility family Bitmask types. Read its declarations together with mandates, preconditions, effects, results, exceptions, complexity, and lifetime rules.
\item[16.3.3.3.4] \textbf{Character sequences.} Specifies text-sequence behavior for Character sequences; establish character type, extent, ownership, null termination, encoding assumptions, and invalidation before use.
\item[16.3.3.3.4.1] \textbf{General.} Frames the terminology, applicability, and relationships needed to interpret General; detailed rules remain in the following subclauses.
\item[16.3.3.3.4.2] \textbf{Byte strings.} Specifies text-sequence behavior for Byte strings; establish character type, extent, ownership, null termination, encoding assumptions, and invalidation before use.
\item[16.3.3.3.4.3] \textbf{Multibyte strings.} Specifies text-sequence behavior for Multibyte strings; establish character type, extent, ownership, null termination, encoding assumptions, and invalidation before use.
\item[16.3.3.3.5] \textbf{Customization Point Object types.} Specifies the facility family Customization Point Object types. Read its declarations together with mandates, preconditions, effects, results, exceptions, complexity, and lifetime rules.
\item[16.3.3.4] \textbf{Functions within classes.} Specifies the facility family Functions within classes. Read its declarations together with mandates, preconditions, effects, results, exceptions, complexity, and lifetime rules.
\item[16.3.3.5] \textbf{Private members.} Specifies the facility family Private members. Read its declarations together with mandates, preconditions, effects, results, exceptions, complexity, and lifetime rules.
\item[16.3.3.6] \textbf{Freestanding items.} Specifies the facility family Freestanding items. Read its declarations together with mandates, preconditions, effects, results, exceptions, complexity, and lifetime rules.
\end{description}
\Needspace{8\baselineskip}
\subsection*{16.4 Library-wide requirements}
\begin{longtable}{@{}>{\RaggedRight\arraybackslash\columncolor{CornellReferenceCueGray}}p{0.28\textwidth}>{\RaggedRight\arraybackslash}p{0.64\textwidth}@{}}
\rowcolor{CornellReferenceNavy}\color{white}\textbf{Cue / recall prompt} & \color{white}\textbf{Notes}\\\endfirsthead
\rowcolor{CornellReferenceNavy}\color{white}\textbf{Cue / recall prompt} & \color{white}\textbf{Notes (continued)}\\\endhead
\arrayrulecolor{CornellReferenceLineGray}
\textbf{What contract governs General?}\\[-1mm]{\scriptsize\CornellClauseRef{16.4.1}} & Frames the terminology, applicability, and relationships needed to interpret General; detailed rules remain in the following subclauses.\\\hline
\textbf{What contract governs Library contents and organization?}\\[-1mm]{\scriptsize\CornellClauseRef{16.4.2}} & Specifies the facility family Library contents and organization. Read its declarations together with mandates, preconditions, effects, results, exceptions, complexity, and lifetime rules.\\\hline
\textbf{What contract governs Using the library?}\\[-1mm]{\scriptsize\CornellClauseRef{16.4.3}} & Specifies the facility family Using the library. Read its declarations together with mandates, preconditions, effects, results, exceptions, complexity, and lifetime rules.\\\hline
\textbf{Which obligations govern Requirements on types and expressions?}\\[-1mm]{\scriptsize\CornellClauseRef{16.4.4}} & States the admissibility and semantic obligations for Requirements on types and expressions. Separate compile-time participation rules from caller preconditions and runtime effects.\\\hline
\textbf{What contract governs Constraints on programs?}\\[-1mm]{\scriptsize\CornellClauseRef{16.4.5}} & States the admissibility and semantic obligations for Constraints on programs. Separate compile-time participation rules from caller preconditions and runtime effects.\\\hline
\textbf{What contract governs Conforming implementations?}\\[-1mm]{\scriptsize\CornellClauseRef{16.4.6}} & Specifies the facility family Conforming implementations. Read its declarations together with mandates, preconditions, effects, results, exceptions, complexity, and lifetime rules.\\\hline
\end{longtable}
\begin{CornellReferenceSummaryBox}Library-wide requirements supplies the requirements portion of this clause. The key review move is to connect its local wording to the clause-wide model, then classify each condition by normative force before relying on it.\end{CornellReferenceSummaryBox}
\Needspace{5\baselineskip}\paragraph{Detailed coverage ledger.}
\begin{description}[style=nextline,leftmargin=2.8cm,font=\normalfont\bfseries\color{CornellReferenceTeal}]
\item[16.4.2.1] \textbf{General.} Frames the terminology, applicability, and relationships needed to interpret General; detailed rules remain in the following subclauses.
\item[16.4.2.2] \textbf{Library contents.} Specifies the facility family Library contents. Read its declarations together with mandates, preconditions, effects, results, exceptions, complexity, and lifetime rules.
\item[16.4.2.3] \textbf{Headers.} Specifies the facility family Headers. Read its declarations together with mandates, preconditions, effects, results, exceptions, complexity, and lifetime rules.
\item[16.4.2.4] \textbf{Modules.} Places Modules within the module ownership and dependency model; visibility, reachability, linkage, and textual preprocessing remain distinct.
\item[16.4.2.5] \textbf{Freestanding implementations.} Specifies the facility family Freestanding implementations. Read its declarations together with mandates, preconditions, effects, results, exceptions, complexity, and lifetime rules.
\item[16.4.3.1] \textbf{Overview.} Frames the terminology, applicability, and relationships needed to interpret Overview; detailed rules remain in the following subclauses.
\item[16.4.3.2] \textbf{Headers.} Specifies the facility family Headers. Read its declarations together with mandates, preconditions, effects, results, exceptions, complexity, and lifetime rules.
\item[16.4.3.3] \textbf{Linkage.} Locates names and entity identity for Linkage; distinguish where a name is visible from whether declarations denote the same entity across contexts.
\item[16.4.4.1] \textbf{General.} Frames the terminology, applicability, and relationships needed to interpret General; detailed rules remain in the following subclauses.
\item[16.4.4.2] \textbf{Template argument requirements.} States the admissibility and semantic obligations for Template argument requirements. Separate compile-time participation rules from caller preconditions and runtime effects.
\item[16.4.4.3] \textbf{Swappable requirements.} Governs state transfer or exchange for Swappable requirements; check self-operation, validity after movement, exception behavior, and invalidation.
\item[16.4.4.4] \textbf{Cpp17NullablePointer requirements.} States the admissibility and semantic obligations for Cpp17NullablePointer requirements. Separate compile-time participation rules from caller preconditions and runtime effects.
\item[16.4.4.5] \textbf{Cpp17Hash requirements.} States the admissibility and semantic obligations for Cpp17Hash requirements. Separate compile-time participation rules from caller preconditions and runtime effects.
\item[16.4.4.6] \textbf{Cpp17Allocator requirements.} States the admissibility and semantic obligations for Cpp17Allocator requirements. Separate compile-time participation rules from caller preconditions and runtime effects.
\item[16.4.4.6.1] \textbf{General.} Frames the terminology, applicability, and relationships needed to interpret General; detailed rules remain in the following subclauses.
\item[16.4.4.6.2] \textbf{Allocator completeness requirements.} States the admissibility and semantic obligations for Allocator completeness requirements. Separate compile-time participation rules from caller preconditions and runtime effects.
\item[16.4.5.1] \textbf{Overview.} Frames the terminology, applicability, and relationships needed to interpret Overview; detailed rules remain in the following subclauses.
\item[16.4.5.2] \textbf{Namespace use.} Specifies the facility family Namespace use. Read its declarations together with mandates, preconditions, effects, results, exceptions, complexity, and lifetime rules.
\item[16.4.5.2.1] \textbf{Namespace std.} Specifies the facility family Namespace std. Read its declarations together with mandates, preconditions, effects, results, exceptions, complexity, and lifetime rules.
\item[16.4.5.2.2] \textbf{Namespace posix.} Specifies the facility family Namespace posix. Read its declarations together with mandates, preconditions, effects, results, exceptions, complexity, and lifetime rules.
\item[16.4.5.2.3] \textbf{Namespaces for future standardization.} Defines asynchronous shared-state behavior for Namespaces for future standardization; establish producer/consumer ownership, readiness, blocking, exception transport, and destruction effects.
\item[16.4.5.3] \textbf{Reserved names.} Specifies the facility family Reserved names. Read its declarations together with mandates, preconditions, effects, results, exceptions, complexity, and lifetime rules.
\item[16.4.5.3.1] \textbf{General.} Frames the terminology, applicability, and relationships needed to interpret General; detailed rules remain in the following subclauses.
\item[16.4.5.3.2] \textbf{Zombie names.} Specifies the facility family Zombie names. Read its declarations together with mandates, preconditions, effects, results, exceptions, complexity, and lifetime rules.
\item[16.4.5.3.3] \textbf{Macro names.} Operates on preprocessing tokens for Macro names; expansion, argument substitution, rescanning, and directive context occur before core-language parsing.
\item[16.4.5.3.4] \textbf{External linkage.} Locates names and entity identity for External linkage; distinguish where a name is visible from whether declarations denote the same entity across contexts.
\item[16.4.5.3.5] \textbf{Types.} Specifies the facility family Types. Read its declarations together with mandates, preconditions, effects, results, exceptions, complexity, and lifetime rules.
\item[16.4.5.3.6] \textbf{User-defined literal suffixes.} Specifies the facility family User-defined literal suffixes. Read its declarations together with mandates, preconditions, effects, results, exceptions, complexity, and lifetime rules.
\item[16.4.5.4] \textbf{Headers.} Specifies the facility family Headers. Read its declarations together with mandates, preconditions, effects, results, exceptions, complexity, and lifetime rules.
\item[16.4.5.5] \textbf{Derived classes.} Specifies the facility family Derived classes. Read its declarations together with mandates, preconditions, effects, results, exceptions, complexity, and lifetime rules.
\item[16.4.5.6] \textbf{Replacement functions.} Specifies the facility family Replacement functions. Read its declarations together with mandates, preconditions, effects, results, exceptions, complexity, and lifetime rules.
\item[16.4.5.7] \textbf{Handler functions.} Specifies the facility family Handler functions. Read its declarations together with mandates, preconditions, effects, results, exceptions, complexity, and lifetime rules.
\item[16.4.5.8] \textbf{Other functions.} Specifies the facility family Other functions. Read its declarations together with mandates, preconditions, effects, results, exceptions, complexity, and lifetime rules.
\item[16.4.5.9] \textbf{Function arguments.} Specifies the facility family Function arguments. Read its declarations together with mandates, preconditions, effects, results, exceptions, complexity, and lifetime rules.
\item[16.4.5.10] \textbf{Library object access.} Specifies the facility family Library object access. Read its declarations together with mandates, preconditions, effects, results, exceptions, complexity, and lifetime rules.
\item[16.4.5.11] \textbf{Semantic requirements.} States the admissibility and semantic obligations for Semantic requirements. Separate compile-time participation rules from caller preconditions and runtime effects.
\item[16.4.6.1] \textbf{Overview.} Frames the terminology, applicability, and relationships needed to interpret Overview; detailed rules remain in the following subclauses.
\item[16.4.6.2] \textbf{Headers.} Specifies the facility family Headers. Read its declarations together with mandates, preconditions, effects, results, exceptions, complexity, and lifetime rules.
\item[16.4.6.3] \textbf{Restrictions on macro definitions.} Operates on preprocessing tokens for Restrictions on macro definitions; expansion, argument substitution, rescanning, and directive context occur before core-language parsing.
\item[16.4.6.4] \textbf{Non-member functions.} Specifies the facility family Non-member functions. Read its declarations together with mandates, preconditions, effects, results, exceptions, complexity, and lifetime rules.
\item[16.4.6.5] \textbf{Member functions.} Specifies the facility family Member functions. Read its declarations together with mandates, preconditions, effects, results, exceptions, complexity, and lifetime rules.
\item[16.4.6.6] \textbf{Friend functions.} Specifies the facility family Friend functions. Read its declarations together with mandates, preconditions, effects, results, exceptions, complexity, and lifetime rules.
\item[16.4.6.7] \textbf{Constexpr functions and constructors.} Specifies how Constexpr functions and constructors establishes object state, including participation constraints, initialization effects, and exceptional exits.
\item[16.4.6.8] \textbf{Requirements for stable algorithms.} States the admissibility and semantic obligations for Requirements for stable algorithms. Separate compile-time participation rules from caller preconditions and runtime effects.
\item[16.4.6.9] \textbf{Reentrancy.} Specifies the facility family Reentrancy. Read its declarations together with mandates, preconditions, effects, results, exceptions, complexity, and lifetime rules.
\item[16.4.6.10] \textbf{Data race avoidance.} Specifies the facility family Data race avoidance. Read its declarations together with mandates, preconditions, effects, results, exceptions, complexity, and lifetime rules.
\item[16.4.6.11] \textbf{Protection within classes.} Specifies the facility family Protection within classes. Read its declarations together with mandates, preconditions, effects, results, exceptions, complexity, and lifetime rules.
\item[16.4.6.12] \textbf{Derived classes.} Specifies the facility family Derived classes. Read its declarations together with mandates, preconditions, effects, results, exceptions, complexity, and lifetime rules.
\item[16.4.6.13] \textbf{Restrictions on exception handling.} Defines the failure model for Restrictions on exception handling, including how an error is represented, reported, propagated, or converted into a diagnostic consequence.
\item[16.4.6.14] \textbf{Value of error codes.} Defines the failure model for Value of error codes, including how an error is represented, reported, propagated, or converted into a diagnostic consequence.
\item[16.4.6.15] \textbf{Moved-from state of library types.} Specifies the facility family Moved-from state of library types. Read its declarations together with mandates, preconditions, effects, results, exceptions, complexity, and lifetime rules.
\end{description}
\section{Language-Lawyer and Library-Usage Pitfalls}
\begin{CornellReferencePitfallBox}\begin{itemize}[label=\textcolor{CornellReferenceAmber}{$\blacktriangleright$}]
\item Reading a synopsis as the complete behavioral contract.
\item Ignoring mandates, preconditions, exception behavior, or complexity requirements.
\item Using an exposition-only name as though it were public API.
\item Assuming invalidation, ownership, or thread-safety properties that are not stated.
\item Treating successful compilation as proof that semantic requirements and preconditions are met.
\end{itemize}\end{CornellReferencePitfallBox}
\section{Programmer and Implementation Checklist}
\begin{CornellReferenceChecklistBox}\begin{itemize}[label=$\square$]
\item Identify the declaring header or module exposure for each facility.
\item Check template arguments and mandates before evaluating an operation.
\item Establish every precondition at the call site.
\item Record effects, returned value, postconditions, and exceptions separately.
\item Audit iterator, reference, pointer, and object lifetime after mutation.
\item Verify stated complexity and synchronization assumptions.
\item For \CornellClauseRef{16.1}, verify general against the precise rule category and all cited dependencies.
\item For \CornellClauseRef{16.2}, verify the c standard library against the precise rule category and all cited dependencies.
\item For \CornellClauseRef{16.3}, verify method of description against the precise rule category and all cited dependencies.
\item For \CornellClauseRef{16.4}, verify library-wide requirements against the precise rule category and all cited dependencies.
\end{itemize}\end{CornellReferenceChecklistBox}
\section{Clause Synthesis}
\begin{CornellReferenceOrientationBox}Defines how library clauses are organized and interpreted, then states requirements on programs and conforming implementations, including reserved names, user specializations, type requirements, moved-from states, reentrancy, and data-race avoidance. The interpretation and conformance framework for every library facility. A sound review distinguishes syntax and participation from runtime preconditions, keeps notes and examples non-mandatory, and follows cross-clause dependencies before making a portability claim.\end{CornellReferenceOrientationBox}
\section{Self-Test Questions}
\begin{enumerate}
\item What role does General play in the clause's overall model?
\item Which normative distinctions must be kept separate when applying The C standard library?
\item How would you audit a program or implementation that depends on Method of description?
\item Compare Library-wide requirements with General; what different question does each answer?
\item What role does General play in the clause's overall model?
\item Which normative distinctions must be kept separate when applying Structure of each clause?
\item How would you audit a program or implementation that depends on Other conventions?
\item Compare General with Library contents and organization; what different question does each answer?
\item What role does Library contents and organization play in the clause's overall model?
\item Which normative distinctions must be kept separate when applying Using the library?
\item How would you audit a program or implementation that depends on Requirements on types and expressions?
\item Compare Constraints on programs with Conforming implementations; what different question does each answer?
\item What role does Conforming implementations play in the clause's overall model?
\item Which normative distinctions must be kept separate when applying mandate?
\item Scenario: a program relies on an unstated implementation behavior while using precondition. What must be checked before calling the program portable?
\item Scenario: a program relies on an unstated implementation behavior while using exposition-only entity. What must be checked before calling the program portable?
\end{enumerate}
\clearpage\section{Answer Key}
\begin{enumerate}
\item General is one part of the clause's library semantics model. Its role is understood by connecting the local rule in 16.4.5.3.1 to the clause orientation and its dependent concepts.
\item Keep formation and well-formedness separate from runtime preconditions and effects. Also distinguish mandatory wording from notes, implementation choice from undefined behavior, and interface declarations from detailed contracts; see 16.2.
\item Audit Method of description by locating the controlling wording in 16.3, listing inputs and type constraints, proving any preconditions, recording effects and failure behavior, and checking lifetime, invalidation, complexity, and synchronization where applicable.
\item Library-wide requirements and General occupy different steps or facility groups. Analyze each under its own subclause, then follow the explicit dependency between them rather than transferring one set of guarantees to the other; begin at 16.4.
\item General is one part of the clause's library semantics model. Its role is understood by connecting the local rule in 16.4.5.3.1 to the clause orientation and its dependent concepts.
\item Keep formation and well-formedness separate from runtime preconditions and effects. Also distinguish mandatory wording from notes, implementation choice from undefined behavior, and interface declarations from detailed contracts; see 16.3.2.
\item Audit Other conventions by locating the controlling wording in 16.3.3, listing inputs and type constraints, proving any preconditions, recording effects and failure behavior, and checking lifetime, invalidation, complexity, and synchronization where applicable.
\item General and Library contents and organization occupy different steps or facility groups. Analyze each under its own subclause, then follow the explicit dependency between them rather than transferring one set of guarantees to the other; begin at 16.4.5.3.1.
\item Library contents and organization is one part of the clause's library semantics model. Its role is understood by connecting the local rule in 16.4.2 to the clause orientation and its dependent concepts.
\item Keep formation and well-formedness separate from runtime preconditions and effects. Also distinguish mandatory wording from notes, implementation choice from undefined behavior, and interface declarations from detailed contracts; see 16.4.3.
\item Audit Requirements on types and expressions by locating the controlling wording in 16.4.4, listing inputs and type constraints, proving any preconditions, recording effects and failure behavior, and checking lifetime, invalidation, complexity, and synchronization where applicable.
\item Constraints on programs and Conforming implementations occupy different steps or facility groups. Analyze each under its own subclause, then follow the explicit dependency between them rather than transferring one set of guarantees to the other; begin at 16.4.5.
\item Conforming implementations is one part of the clause's library semantics model. Its role is understood by connecting the local rule in 16.4.6 to the clause orientation and its dependent concepts.
\item Keep formation and well-formedness separate from runtime preconditions and effects. Also distinguish mandatory wording from notes, implementation choice from undefined behavior, and interface declarations from detailed contracts; see 16.
\item First identify the exact requirement in 16, then classify the relied-on behavior as required, unspecified, implementation-defined, conditionally supported, or undefined. Portability requires satisfying the program obligations and avoiding dependence on a choice not guaranteed in the target environments.
\item First identify the exact requirement in 16, then classify the relied-on behavior as required, unspecified, implementation-defined, conditionally supported, or undefined. Portability requires satisfying the program obligations and avoiding dependence on a choice not guaranteed in the target environments.
\end{enumerate}
\section{Key-Term Recap}
\begin{longtable}{@{}>{\RaggedRight\arraybackslash}p{0.28\textwidth}>{\RaggedRight\arraybackslash}p{0.64\textwidth}@{}}\toprule\textbf{Term} & \textbf{Working meaning}\\\midrule\endfirsthead\toprule\textbf{Term} & \textbf{Working meaning (continued)}\\\midrule\endhead\bottomrule\endfoot
mandate & A declaration-level requirement whose failure affects participation or well-formedness as specified.\\
precondition & A caller obligation required before an operation is evaluated.\\
exposition-only entity & A specification device that is not public program API.\\
freestanding item & A library entity required for the freestanding environment.\\
moved-from state & A valid but otherwise unspecified library-object state unless a stronger contract is given.\\
reserved name & An identifier or name pattern withheld from ordinary program definition or use.\\
\end{longtable}
\CornellTakeaway{Library wording is a structured contract whose labels, program constraints, and implementation obligations must be read together.}
\end{document}
